\subsection{Collider}

Collider je komponenta v Unity, která umožňuje přidat kolize na objekty uvnitř hry. Díky colliderům se dá poté například určovat ničení herních objektů nebo zabíjení nepřátel. Také lze colliderům nastavit funkci \textit{Trigger}, která se používá například při aktivování jiných scén ve hře nebo při aktivaci interakcí. \textit{Trigger collider} nezastaví hráče při kolizi, ale detekuje průchod a následně spustí požadovanou akci nebo kód.

\subsection{Rigidbody}

Rigidbody je komponenta v Unity, která umožňuje herním objektům používat fyzikální engine hry. Díky manipulaci herních objektů upravováním parametrů v Rigidbody je možné dělat realistické kolize, upravovat gravitaci nebo celkově nastavit chování herních objektů.

\subsection{Player controller}

Označení \textit{Player controller} se nejčastěji používá pro hlavní skript, který se stará o ovládání hráče. Do tohoto skriptu se nejčastěji přidávají funkce pro chůzi, skákání nebo jiné pokročilejší mechaniky, které jsou specifické pro danou hru.

\subsection{Input action}

Input action je speciální komponent v Unity, který řeší input z klávesnice, myši nebo třeba controlleru pro konzole. Dříve se programoval input za pomoci \texttt{Input.GetKey(KeyCode.\_)}, což nebylo vyloženě špatné, ovšem input action je funkční zjednodušení tohoto systému.

\subsection{Sprite}

Sprite je grafický asset vytvořený ve specializovaném programu. Jedná se o komponenty, které člověk uvidí ve výsledném produktu.

\subsection{Animace}

Animace je složení několika spritů za sebe, které po aktivaci navozují pocit pohybu. Lze u nich upravovat rychlost a délku.

\subsection{Animator controller}

Animator controller neboli zkráceně \textit{Animator} se využívá ke spojování animací ve hře a vytváření hierarchie mezi nimi. Ve spolupráci s kódem lze nastavovat spouštění specifických animací při specifických příležitostech.

\subsection{Tilemap}

Tilemap je metoda tvorby herního prostředí pomocí pravidelné mřížky složené z grafických dlaždic (tiles). Každá dlaždice představuje samostatný grafický prvek, který je umístěn do určité buňky mřížky. Tento systém se používá především ve 2D hrách, kde umožňuje efektivní a přehlednou tvorbu rozsáhlých herních map. Výhodou tilemap systému je možnost opakovaného použití grafických prvků, snížení paměťové náročnosti a snadná úprava úrovní.

\subsection{Tileset}

Tileset je soubor grafických dlaždic určených k sestavení herního prostředí. Obvykle je tvořen jedním obrázkem (sprite sheet), který obsahuje více dlaždic stejné velikosti uspořádaných vedle sebe. Použití tilesetu umožňuje efektivní práci s grafikou, protože jednotlivé prvky lze opakovaně využívat při tvorbě mapy.

\subsection{Game feel}

\textit{Game feel} neboli herní pocit je termín, který popisuje, jak hra působí na hráče z hlediska ovládání a reakcí. Zahrnuje faktory jako je responsivita ovládání, plynulost animací, zpětná vazba od hry a celkový dojem z interakce s herním světem. Kvalitní herní pocit je to, co odlišuje průměrné hry od těch výborných.

\subsection{Frame}

Frame neboli snímek je jednotka času v herním cyklu. Hry typicky běží při 30 až 60 snímcích za sekundu, což znamená, že každou sekundu se herní logika a grafika aktualizují 30 až 60 krát. Pochopení framu je důležité pro implementaci herních mechanik, protože mnoho systémů (jako coyote time nebo jump buffer) pracuje právě s časovým intervalem měřeným ve snímcích.

\subsection{Layer}

Layer neboli vrstva je systém v Unity pro organizaci a filtrování objektů ve scéně. Umožňuje například určit, které objekty mají detekovat kolize s jinými objekty. V kontextu platformovek se často hovoří o \textit{ground layer}, což je vrstva obsahující všechny platformy a zem, na které může hráč stát.

\subsection{Prefab}

Prefab (prefabricated object) je předem vytvořený herní objekt uložený jako asset v projektu. Umožňuje vytvořit instanci (kopii) objektu přímo v kódu pomocí metody \newline \texttt{Instantiate}. Prefaby se často používají pro projektil, nepřátele nebo jiné opakovaně vytvářené objekty ve hře.

\subsection{ScriptableObject}

ScriptableObject je speciální datový typ v Unity určený pro ukládání konfiguračních dat a parametrů. Na rozdíl od \texttt{MonoBehaviour} není přidělen k žádnému hernímu objektu ve scéně a může existovat jako samostatný asset. Používá se typicky pro nastavení herních parametrů, které je možné měnit v editoru bez nutnosti upravovat kód.

\subsection{Callback}

Callback je mechanismus v programování, kdy se jako parametr předá jiná funkce, která se má zavolat po dokončení určité operace. V kontextu Unity input systému se callbacky používají k reagování na události, jako je stisknutí nebo uvolnění klávesy. Když hráč stiskne tlačítko, systém automaticky zavolá předem definovanou metodu (callback).